\chapter{Forschungsstand}\label{ch:stand_der_forschung}

Die aktuelle Forschung in diesem Gebiet verbindet Sprachmodelle mit 3D-Repräsentationen und zeigt, dass LLM-basierte Systeme zunehmend in der Lage sind, 3D-Szenen nicht nur mit Sprache zu beschreiben, sondern semantisch sowie räumlich zu interpretieren.
Mit der Weiterentwicklung von LLMs konnte deren Integration mit räumlichen Daten große Fortschritte erzielen und bietet so neue Möglichkeiten wie Szenenverständnis, Planung und Generierung von 3D-Welten oder Navigation in diesen.
Spätestens seit 2023 hat sich die Verbindung zwischen LLMs und 3D-Repräsentationen zu einer eigenständigen Forschungslinie entwickelt, welche 3D-Szenen nicht nur beschreiben, sondern auch für neue Methoden wie das \textit{Grounding}, Fragebeantwortung und Planungsaufgaben nutzen können~\cite{maWhenLLMsStep2025,hong3DLLMInjecting3D2023}.
Trotz dieser schnellen Entwicklung bleiben Datenknappheit, Halluzinationsrobustheit und räumliches Verständnis (Grounding) zentrale Herausforderungen aktueller Systeme~\cite{pengUnderstandingEvaluatingHallucinations2025, maWhenLLMsStep2025}.

\paragraph{3D-LLMs für 3D-Verstehen.}

In Arbeiten wie 3D-LLM und PointLLM können multimodale LLMs 3D-Daten direkt verarbeiten und nehmen dafür beispielsweise Punktwolken als Eingabe entgegen, wobei PointLLM eher objektzentriert ist. Dadurch können Aufgaben wie \textit{3D-Captioning}, 3D-Fragebeantwortung, Grounding, Dialog oder Navigation bearbeitet werden und relevante komplexere Konzepte wie räumliche Beziehungen zwischen Objekten, physikalische Eigenschaften, Raumlayouts oder Interaktionsmöglichkeiten (\textit{Affordances}) verstehen~\cite{hong3DLLMInjecting3D2023,xuPointLLMEmpoweringLarge2025}.

\paragraph{LLM-gestützte Pipelines zur Szenengenerierung und -bearbeitung.}

Parallel dazu wurden LLM-gestützte Pipelines entwickelt, in denen LLMs nicht direkt 3D-Daten verarbeiten.
Stattdessen fungieren sie als Planungs- und Steuerungseinheit, die für die Generierung und Bearbeitung ganzer 3D-Szenen genutzt werden.
Das Framework LLMR nutzt für die Generierung von interaktiven Szenen einen modularen Workflow, darunter die Module Planner, Scene Analyzer, Skill Library, Builder und Inspector. Es kann Szenenkontext analysieren und generiert anschließend Unity-kompatiblen C\#-Code, welcher durch ein Inspektionsmodul iterativ geprüft wird. Als Eingabe arbeitet das System mit einer strukturierten, für das LLM konsumierbaren JSON-Repräsentation des Szenenkontexts~\cite{delatorreLLMRRealtimePrompting2024}.
HOLODECK verfolgt einen stärker szenenzentrierten Ansatz und zerlegt die Generierung einer Szene in einzelne Schritte für Grundrisse, Materialien, Fenster und Türen, Objektauswahl sowie räumliche Anordnung.
Es dient der automatisierten Erstellung von 3D-Simulationsumgebungen im Bereich der verkörperten KI (\textit{Embodied-AI}) und platziert durch ein Constraint-basiertes Verfahren Objekte räumlich plausibel im Raum~\cite{yangHolodeckLanguageGuided2024}.

\section{Artikulation statischer 3D-Objekte}
%TODO Artikulierung ersetzen

Für die Transformation statischer in interaktive Objekte sind objekt- und szenenzentrierte Ansätze besonders relevant.
ArtLLM wandelt eine Punktwolke in eine Abfolge von \textit{Tokens} um, die vom Modell nacheinander generiert werden.
Darüber hinaus können Assets auch als Bild oder Text eingegeben werden. Wichtig ist dabei, dass diese nicht direkt in das Kernmodell eingehen, sondern vorher mit Hilfe externer Generierungsmodelle in initiale Meshes und dann in Punktwolken umgewandelt werden.
Diese Struktur definiert die Positionen der einzelnen Teile, die mechanischen Gelenke sowie deren Hierarchie.
Um Positionen und Winkel anzugeben, arbeitet es mit festen Rasterstufen statt Gleitkommazahlen.
Nachdem detailliertere Geometrie über ein Zusatzmodul hinzugefügt wurde, wird diese kinematische Struktur als URDF-Datei (\textit{Unified Robotics Description Format}) exportiert und kann so in Simulationen und virtuellen Umgebungen verwendet werden~\cite{wangArtLLMGeneratingArticulated2026}.

ArtiWorld erweitert diesen Ansatz von der Objekt- auf die Szenenebene.
Das System nimmt eine textuelle Szenenbeschreibung entgegen und identifiziert auf dieser Grundlage artikulierbare Objekte.
Anschließend erzeugt es strukturelle Beschreibungen der einzelnen Beziehungen und anschließend ebenfalls vollständige URDF-Dateien, welche an die ursprünglichen Koordinaten des starren 3D-Assets angepasst werden~\cite{yangArtiWorldLLMDrivenArticulation2025}.
Die Geometrie des Originals wird dabei nicht ersetzt, sondern um Interaktivität erweitert.
Insbesondere diese Arbeit zeigt, dass strukturelle Zwischenrepräsentationen das Reasoning verbessern.


\section{Bildbasierte Eingaben}

Eingaben basierend auf Bildern gewinnen in artikulationsbezogenen Pipelines an Bedeutung.
Einzelbilder werden aktuell in mehreren Arbeiten genutzt, um simulierbare Szenen oder artikulierbare Objekte zu generieren, wie etwa in URDFormer oder ArtLLM~\cite{chenURDFormerPipelineConstructing2024,wangArtLLMGeneratingArticulated2026}.
Dabei können Bilder sowohl isoliert, wie etwa in URDFormer, als auch zur ergänzenden Eingabe, wie in ArtLLM, verwendet werden, da diese semantische Hinweise und Informationen zu potenziellen Interaktionsmöglichkeiten bieten.
Für die vorliegende Arbeit ist dieser Befund besonders relevant.
Er zeigt die Kombination aus expliziten Raumdaten in Form von strukturellen Repräsentationen wie Punktwolken oder Szenengraphen mit Transformationen und Hierarchien einerseits sowie Bildern als visuellem Grounding andererseits.

\section{Forschungslücke}\label{sec:forschungslücke}
% und Relevanz der/für die vorliegende/n Arbeit

Die in diesem Kapitel betrachteten Arbeiten zeigen, dass LLMs und MLLMs zunehmend in 3D-bezogenen Aufgabenbereichen eingesetzt werden und viele relevante Teilansätze wie 3D-Szenenverständnis, 3D-Fragebeantwortung, 3D-Grounding, Interaktions- und Aufgabenplanung, textbasierte 3D-Verarbeitung, Szenengenerierung sowie artikulationsbezogene Objektmodellierung existieren.
So adressiert 3D-LLM \ua~3D-Captioning, 3D-Fragebeantwortung, 3D-Grounding, Dialog und Navigation~\cite{hong3DLLMInjecting3D2023}.
LLMR nutzt LLMs zur Erstellung und Modifikation interaktiver 3D-Welten zusammen mit Modulen wie Planner, Scene Analyzer, Builder und Inspector~\cite{delatorreLLMRRealtimePrompting2024}.
Holodeck erzeugt 3D-Umgebungen aus Sprache und nutzt ebenfalls LLMs \ua für Objektwahl und räumliche Layout-Constraints~\cite{yangHolodeckLanguageGuided2024}.
ArtiWorld und URDFormer adressieren artikulationsbezogene \bzw URDF-basierte Simulation und Kinematik~\cite{yangArtiWorldLLMDrivenArticulation2025,chenURDFormerPipelineConstructing2024}.
Diese Ansätze adressieren jedoch meist andere Zielartefakte als die vorliegende Arbeit und erzeugen häufig neue 3D-Geometrien, bearbeitete 3D-Assets, vollständige 3D-Szenen, Unity-Code, URDF-Dateien oder robotische Simulationsumgebungen.
Die vorliegende Arbeit ist im Gegensatz dazu nicht primär als Verfahren zur Geometrie- oder Szenengenerierung einzuordnen, sondern als semantisch-verhaltensbezogener Ansatz zur Ableitung von Interaktionsspezifikationen aus vorhandenen 3D-Szenen.

Die Forschungslücke ergibt sich damit aus der Verbindung mehrerer Aspekte.
Es sollen vorhandene statische 3D-Objekte erhalten bleiben, relevante Objekte und Teilobjekte semantisch interpretiert werden und gewünschtes Verhalten soll natürlichsprachlich beschrieben werden können.
Darüber hinaus ergibt sich aus den Randbedingungen des Vivian Frameworks, dass eine maschinenlesbare Interaktionslogik entstehen soll, welche kompatibel mit dem Vivian-Framework ist.
Da LLM-gestützte Generierungsprozesse fehlerhafte Objektverweise, inkonsistente Strukturen oder unzulässige Spezifikationselemente erzeugen können, sind zudem Validierungs- und Korrekturmaßnahmen erforderlich.
Genau diese Kombination wird im derzeitigen Stand der Forschung nicht direkt adressiert.
Die Problemstellung in Kapitel~\ref{ch:einführung} nennt bereits, dass statische 3D-Objekte zwar Geometrie, Materialien und Transformationen besitzen, es aber an Wissen über bedienbare Elemente, Funktionen und Reaktionen auf Nutzer*innenhandlungen fehlt.
Dort wird außerdem bereits ein niedrigschwelliger und kontrollierbarer Ansatz gefordert, welcher vorhandene statische 3D-Objekte in maschinenlesbare, validierbare und zustandsbasierte Prototypen überführt, während die Geometrie erhalten bleiben soll.

Aus dieser Forschungslücke sowie den Forschungsfragen F1 bis F4 (siehe~\ref{sec:forschungsfragen_zielsetzung}) folgt die Notwendigkeit für ein System, welches Szeneninformationen analysiert, Nutzer*innenbeschreibungen von Interaktionen verarbeitet und daraus relevante Objekte, Teilobjekte sowie Rollen ableitet, Vivian-kompatible Interaktionsspezifikationen erzeugt und Fehler prüft sowie gezielt korrigiert.

% Aufbauend auf den in Kapitel~\ref{ch:einführung} aufgeführten Problemstellungen und Forschungsfragen, dem Stand der Forschung sowie den in~\ref{sec:vivian_framework} beschriebenen technischen Randbedingungen des Vivian-Frameworks werden zunächst zentrale Problemfelder identifiziert, die durch das in dieser Arbeit implementierte System adressiert werden sollen.
% Wie in der Problemstellung gezeigt, besitzen statische 3D-Objekte an sich zwar geometrische und visuelle Informationen, aber keine explizite Interaktionssemantik.
% Gleichzeitig erschwert manuelles Vorgehen in der Erstellung interaktiver Prototypen nicht-technischen Nutzer*innen den Zugang zu dieser.
% Vivian setzt außerdem voraus, dass Interaktionsverhalten in eine formal korrekte, maschinenlesbare Spezifikation aus Interaktionselementen, Visualisierungselementen, Zuständen und Übergängen überführt wird.
% Da LLM-gestützte Verfahren fehleranfällig sein können~\cite{zhaoSurveyLargeLanguage2026a}, müssen syntaktische, semantische und referenzielle Fehler zudem bereits vor der Laufzeit erkannt und korrigiert werden.

% Aus diesen Überlegungen ergeben sich die Problemfelder P1-P3.
% Die Anforderungen an die in dieser Arbeit implementierte Pipeline werden im Folgenden nicht empirisch durch eine Nutzerstudie erhoben, sondern aus diesen Problemfeldern, den Forschungsfragen und den Vivian-spezifischen Rahmenbedingungen abgeleitet.

% \textbf{P1: Hohe Einstiegshürde für nicht-technische Nutzer*innen.} Die Erstellung interaktiver Prototypen setzt oft Kenntnisse in der Programmierung oder 3D-Modellierung voraus. LLMR zeigt, dass natürliche Sprache zur Erstellung und Modifikation von interaktiven 3D-Szenen genutzt werden kann, es aber zudem spezialisierter Module wie Scene Analyzer, Planner, Builder und Inspector bedarf, da ein einfaches LLM für die interaktive Szenenerstellung nicht ausreicht~\cite{delatorreLLMRRealtimePrompting2024}.

% \textbf{P2: Semantische Adressierbarkeit vorhandener 3D-Objekte und Teilobjekte.} Statische 3D-Objekte enthalten meist primär geometrische, visuelle und transformatorische Informationen und keine darüber, welche Bestandteile bedienbar sind oder welche Rolle sie in einer Interaktion besitzen. Ohne diese semantische Adressierbarkeit kann eine spätere Interaktionsbeschreibung nicht mit der vorhandenen Geometrie verknüpft werden.
% Klassische LLMs und Vision Language Models (VLMs) allein sind nicht ausreichend in der dreidimensionalen Welt verankert, um statische 3D-Objekte in einer virtuellen Welt mit solchen Informationen anzureichern~\cite{hong3DLLMInjecting3D2023}.

% \textbf{P3: Fehlererkennung und Kontrollierbarkeit LLM-gestützter Generierungsprozesse.}
% LLM-gestützte Generierungsprozesse können fehlerhafte Ausgaben erzeugen, wie nicht existierende Objektverweise oder falsche Beziehungen zwischen Objekten. 3D-LLMs können \zB erheblich von Halluzinationen betroffen sein~\cite{pengUnderstandingEvaluatingHallucinations2025}. LLMR adressiert ähnliche Probleme durch einen modularen Aufbau -- ein Inspector-Modul kontrolliert generierten Code vor der Ausführung auf mögliche Kompilierungs- und Laufzeitfehler~\cite{delatorreLLMRRealtimePrompting2024}. Daher nennt die Problemstellung dieser Arbeit explizit die Notwendigkeit interner Prüf- und Korrekturmaßnahmen, um Fehler frühzeitig sowie bereits im Prozess erkennen und beheben zu können.